\problem{P20}{每轮两次缩小的越界时间}{指数等待与奇偶分解}
每轮把当前长度乘以两个新的独立 $\operatorname{Unif}(0,1)$ 变量，所有变量相互独立。从 $R_0=1$ 出发，令 $R_n=\prod_{i=1}^nU_iV_i$，在首次 $R_n<x$ 时停止，$0<x<1$，次数为 $N$。求 $N$ 的分布、均值，以及终点比例 $Y=R_N/x$ 的密度。
\solution 记 $t=-\log x$。所有负对数增量都是参数为 $1$ 的独立指数变量。把每轮的两个增量依次展开，令 $K$ 为到时刻 $t$ 已完成的指数等待次数。对 $k\ge1$，以第 $k$ 次到达时刻 $s$ 积分，
\[
\Prob(K=k)=\int_0^t\frac{s^{k-1}e^{-s}}{(k-1)!}e^{-(t-s)}\dd s
=e^{-t}\frac{t^k}{k!};
\]
$k=0$ 时概率为 $e^{-t}$，故 $K\sim\operatorname{Pois}(t)$。

只有每两个增量结束后才检查长度，所以 $N=\lfloor K/2\rfloor+1$，即
\[
\boxed{\Prob(N=n)=e^{-t}\left[\frac{t^{2n-2}}{(2n-2)!}
+\frac{t^{2n-1}}{(2n-1)!}\right],\quad n\ge1.}
\]
由泊松生成函数，$\Prob(K\text{ 为奇数})=(1-e^{-2t})/2$。于是
\[
\boxed{\E N=1+\frac{\E K-\Prob(K\text{ 为奇数})}{2}
=\frac{3+2t+e^{-2t}}4.}
\]
由指数分布的无记忆性，从 $t$ 到下一次检查的距离 $O=-\log Y$，在 $K$ 为奇数时是一个指数等待，在 $K$ 为偶数时是两个指数等待之和。记 $w_o=(1-e^{-2t})/2$、$w_e=(1+e^{-2t})/2$，则
\[
f_O(o)=e^{-o}(w_o+w_eo),\quad o>0.
\]
作 $y=e^{-o}$ 的变量变换，得到
\[
\boxed{f_Y(y)=w_o-w_e\log y,\quad0<y<1.}
\]
\insight 批量检查改变了越界后的剩余等待。是否已经走完半轮决定还需一个还是两个指数增量，因此终点比例不再均匀。
